% Reference entries, I (AN-11).  Law 1: os/iword/core/edit.c iw_ed_indent
% ([->Indent] C2: to the next tab stop; the rest of the paragraph wraps at
% that stop; WP's C6 kept before the paragraph's [HRt] by iw_ed_fixup),
% spec/iword/wp5_format.h "tabs, indents and hard pages", face_a.c (F4).
% Shots: ../wordperfect51-decomp mint MINT09, MINT10.
% GATES: test-pubs-iword test-iword-codes test-iword-wp5-roundtrip test-iword-face test-flair-iword-block
\Gates{test-pubs-iword test-iword-codes test-iword-wp5-roundtrip test-iword-face test-flair-iword-block}
\begin{Feature}{Indent}
Indent moves the left edge of a whole paragraph in to the next tab stop.
The first line starts at the stop, and every line after it wraps back to
the same stop, not to the left margin.  Use it for numbered lists, quoted
passages and notes.

To indent a paragraph,
\begin{WPSteps}
\item Move the cursor to the start of the paragraph, or to where you want
the indented text to start (after a number, for example).
\item Press \FK{\rarr Indent}{F4}.  The cursor moves to the next tab stop.
\item Type the paragraph.  Press \K{Enter} at its end; the indent ends with
the paragraph.
\end{WPSteps}
\Documents{wpkey}{F4}

\BeginReplay{ref-indent}
For example, type \TypeRun{1.}, press \Press{\rarr Indent}{F4} and type
\TypeRun{Every TPS report must carry the new cover sheet, which is available from the supply room.}
\WordScreen[rows=2]{ref-indent}{An indented paragraph after a number}
\EndReplay

\Notes
\Topic{Code}
Indent puts the code \Code{[\rarr Indent]} in the document.  Delete it to
undo the indent; InitechWord deletes it without asking.

\Topic{Indenting Farther}
Press Indent again to indent to the next tab stop after that: two Indents at
the start of a paragraph indent it one inch.

\Topic{Tab Stops}
The tab stops are every half inch from the left margin.  If you change the
margins, the stops move with the left margin (see \InRef{Tab}).

\SeeAlso{Margins, Left and Right; Tab}
\end{Feature}

\begin{Feature}{Insert}
Everything you type goes into the document at the cursor, in front of the
character the cursor is on.  The text after the cursor moves along to make
room, and the paragraph is wrapped again.  Nothing you type replaces what is
already there; to replace text, delete it and type the new text (see
\InRef{Delete}).

\BeginReplay{ref-insert}
For example, type \TypeRun{Old text}, press \PressKey{Home}, \PressKey{Left}
and type \TypeRun{New } (with the space):
\WordScreen[rows=1]{ref-insert}{Text typed at the start of a line goes in front of what is there}
\EndReplay

\Notes
\Topic{Inserting Codes}
Codes go in at the cursor in the same way: a hard return, a tab, the start
of bold.

\Topic{The Ins Key}
The Ins key does nothing in InitechWord: typing always inserts.

\SeeAlso{Delete; Word Wrap}
\end{Feature}

\begin{Feature}{Italics}
Italic text slants to the right.  It is one of the appearances on the Font
key's Appearance menu, Italc (4); InitechWord draws italic text in italic on
the editing screen.

To type italic text,
\begin{WPSteps}
\item Press \FK{Font}{Ctrl-F8}, select \Ch{Appearance}{2}, then select
\Ch{Italc}{4}.
\item Type the text.
\item Press \KC{Ctrl-F8}{2,4} again to move past the end of the italic.
\end{WPSteps}
The codes are \Code{[ITALC]} and \Code{[italc]}.  To make text you have
already typed italic, block it, press \FK{Font}{Ctrl-F8}, select
\Ch{Appearance}{2} and \Ch{Italc}{4} (see \InRef{Block, Attributes}).

\SeeAlso{Appearance; Block, Attributes; Font}
\end{Feature}
